\documentclass{article}

\usepackage[letterpaper, margin=1in]{geometry}
\usepackage{amsmath,amssymb}
\usepackage{hyperref}
\usepackage{xcolor}

\title{Q34 Definitions and Rules}
\author{Cal 2 Support Notes}
\date{Feb 2026}

\begin{document}
\maketitle

This reference sheet explains key terms and rules used in Question 34.

\section*{Trigonometric substitution}
A method for integrals involving radicals. Common patterns:
\[
\sqrt{a^2-x^2}:\ x=a\sin\theta,\qquad
\sqrt{a^2+x^2}:\ x=a\tan\theta,\qquad
\sqrt{x^2-a^2}:\ x=a\sec\theta.
\]
\textbf{Example:} if $\sqrt{4-x^2}$ appears, use $a=2$ so $x=2\sin\theta$.

\section*{Denominator}
In a fraction $\frac{p(x)}{q(x)}$, the denominator is the bottom expression $q(x)$.
In Q34, the denominator is $(4-x^2)^{3/2}$.

\section*{Fractional exponents}
\[
a^{m/n}=\sqrt[n]{a^m}=\left(a^{1/n}\right)^m.
\]
\textbf{Example:}
\[
(4-x^2)^{3/2}=\left((4-x^2)^{1/2}\right)^3=\left(\sqrt{4-x^2}\right)^3.
\]

\section*{Notation: $\sin^2\theta$}
$\sin^2\theta$ means $(\sin\theta)^2$. This is standard shorthand.

\section*{Pythagorean identity}
\[
\sin^2\theta+\cos^2\theta=1
\quad\Longleftrightarrow\quad
1-\sin^2\theta=\cos^2\theta.
\]

\section*{Rule: $\sqrt{ab}=\sqrt a\sqrt b$}
For nonnegative $a,b$:
\[
\sqrt{ab}=(ab)^{1/2}=a^{1/2}b^{1/2}=\sqrt a\sqrt b.
\]

\section*{Inverse sine $\sin^{-1}(x)$}
$\sin^{-1}(x)$ (also written $\arcsin(x)$) returns the angle whose sine is $x$.
\textbf{Example:} if $\sin\theta=\frac{x}{2}$, then
\[
\theta=\sin^{-1}\!\left(\frac{x}{2}\right).
\]

\end{document}
